%%%%%%%%%%%%%%%%
%%% lev 17 %%%
%%%%%%%%%%%%%%%%

\Chapter{17}

\Verse{1}
{
	\hP{וַיְדַבֵּר יְהֹוָה אֶל מֹשֶׁה לֵּאמֹר׃}
}
{
	\eP{17 And YHWH spoke to Moses, saying,}
}
{
}

\Verse{2}
{
	\hP{דַּבֵּר אֶל אַהֲרֹן וְאֶל בָּנָיו וְאֶל כּל בְּנֵי יִשְׂרָאֵל וְאָמַרְתָּ אֲלֵיהֶם זֶה הַדָּבָר אֲשֶׁר צִוָּה יְהֹוָה לֵאמֹר׃}
}
{
	\eP{
		“Speak to Aaron and to his sons and to all the children of
		Israel and say to them: This is the thing that YHWH has
		commanded, saying:
	}
}
{
}

\Verse{3}
{
	\hP{אִישׁ אִישׁ מִבֵּית יִשְׂרָאֵל אֲשֶׁר יִשְׁחַט שׁוֹר אוֹ כֶשֶׂב אוֹ עֵז בַּמַּחֲנֶה אוֹ אֲשֶׁר יִשְׁחַט מִחוּץ לַמַּחֲנֶה׃}
}
{
	\eP{
		Any man from the house of Israel who slaughters an ox or a
		sheep or a goat in the camp or who slaughters outside the
		camp
	}
}
{
%	\footnote{
%		Any man … who slaughters. Most sacrifices involve the ritual slaughter
%		of an animal that is to be used for food. In Leviticus, YHWH forbids the
%		slaughter of animals outside of a formal ritual setting. For Leviticus
%		there is no such occupation as butcher. If one wants to eat meat, one
%		must bring the animal to the priest, who sacrifices it in the prescribed
%		manner at the altar, retains a portion of it for the sustenance of the
%		priesthood, and gives back the rest for consumption. The priests eat
%		their portion according to prescription as well: specifically, priests
%		eat meat with unleavened bread. As with the laws of Exodus, the reasons
%		for requiring ritual slaughter are not stated explicitly, but there
%		appears to be more at stake than the support of the priesthood. In the
%		initial state of creation in Genesis, humans are vegetarian. Genesis
%		depicts animals, but not plants, as sharing with humans the quality of
%		being living beings (nepeš ayyh,Gen 1:20,21,24,30; 2:7,19). Following
%		the flood, animals are permitted as food. Now the sacrificial laws of
%		Leviticus recognize that this involves the taking of life, and they
%		decree that life cannot be taken as a secular act of slaughter. It must
%		be done in a setting that reflects sensitivity to what is at stake. An
%		indicator that this is the orientation of the Levitical law is that it
%		permits one to hunt game animals such as deer, which one cannot easily
%		lead to an altar, but then one must perform the act of pouring off the
%		blood of the animal because, it says, “the flesh’s life (nepeš) is in
%		the blood” (17:11). At several places in my comments on Leviticus, I
%		have noted that we have lost the feeling of ritual as it was known in
%		the biblical world. Here I would add that if we still had a greater
%		feeling for ritual we would have more of a feeling for the life of an
%		animal. For most humans in this age, we are farther removed from the
%		life and death of the animals (and plants!) that we eat than in any
%		previous generation. We buy meat in sealed packages, so that it no
%		longer resembles the animal from which it came. We have lost the
%		awareness of how this food came to be on our table. Part of life is the
%		consumption of life-forms by other life-forms. But whether one is
%		vegetarian or carnivorous, one will enrich one’s life—and understand
%		better one’s place in the world—if one connects a ritual with that
%		consumption. If it can no longer be sacrifice, then perhaps it can be
%		saying grace, or a blessing, or occasionally visiting a farm or ranch
%		and observing the animals and plants: watching, touching, listening to
%		them.
%	}
}

\Verse{4}
{
	\hP{וְאֶל פֶּתַח אֹהֶל מוֹעֵד לֹא הֱבִיאוֹ לְהַקְרִיב קרְבָּן לַיהֹוָה לִפְנֵי מִשְׁכַּן יְהֹוָה דָּם יֵחָשֵׁב לָאִישׁ הַהוּא דָּם שָׁפָךְ וְנִכְרַת הָאִישׁ הַהוּא מִקֶּרֶב עַמּוֹ׃}
}
{
	\eP{
		and has not brought it to the entrance of the Tent of
		Meeting to bring forward an offering to YHWH in front of
		YHWH’s Tabernacle: blood will be counted to that man. He
		has spilled blood. And that man will be cut off from among
		his people.
	}
}
{
%	\footnote{
%		Tent of Meeting. There is a long additional line here in the Septuagint,
%		Qumran, and Samaritan texts, which has been lost in the Masoretic Text.
%		The verse reads: “and has not brought it to the entrance of the Tent of
%		Meeting to make it a burnt offering or peace offering to YHWH to be
%		acceptable for you as a pleasant smell, and he slaughters it outside and
%		has not brought it to the entrance of the Tent of Meeting …” The eye of
%		the scribe writing the Masoretic Text apparently skipped the line
%		because his eye jumped from one occurrence of the phrase “brought it to
%		the entrance of the Tent of Meeting” to the next. This shows how much of
%		a text can be changed by scribal error and is a warning to take account
%		of all the available versions of the text. (It is also a reminder that
%		systems of seeking patterns or codes by counting the text’s letters are
%		erroneous. See the comment on Exod 1:5.) Footnote 17:4. blood will be
%		counted to that man. He has spilled blood. This person has slaughtered
%		an animal, but the text speaks of it as an equivalent of murdering a
%		human! Why does this act merit such extraordinary severity? Because it
%		merges two of the most essential commandments of the Torah. First, it
%		involves taking a life without respect for life’s sanctity—as discussed
%		in the comment on 17:3. And, second, it violates the law of
%		centralization of worship, which is declared for the first time in this
%		chapter. According to the law, one can sacrifice only at one place on
%		earth: at the altar at the entrance of the Tent of Meeting. If one wants
%		to eat lamb or beef, one must bring one’s sheep or cow to that central
%		altar. In the wilderness it is at the Tent. Later, in the land, it is at
%		the Temple. (The Tent of Meeting is inside the Temple. See the comment
%		on Exod 26:30.) But the all-important point is centralization of
%		worship. As I commented earlier (Exod 26:6), centralization and
%		monotheism go together: one God, one altar. Each town in Israel does not
%		have its own altar or its own Tabernacle. This is explicitly forbidden
%		here. In the books of Kings, it is the standard for judging every king
%		of Israel and Judah: did he allow “high places” (Hebrew bmôt) outside of
%		the Jerusalem Temple or not? The only two kings to get nearly perfect
%		ratings (Hezekiah and Josiah) are the two kings who most zealously
%		enforce centralization. Today, one does not see a multiplicity of
%		synagogues, churches, or mosques and assume that this means that a
%		multiplicity of gods is being worshiped. But in the biblical world, when
%		monotheism is a new idea, the principle of one God–one Temple–one altar
%		is presented as supremely important. Footnote 17:4. cut off from among
%		his people. This may mean to die without heirs or not to be buried in
%		the ancestral tomb or not to join one’s ancestors in an afterlife. See
%		the comment on Exod 30:33.
%	}
}

\Verse{5}
{
	\hP{לְמַעַן אֲשֶׁר יָבִיאוּ בְּנֵי יִשְׂרָאֵל אֶת זִבְחֵיהֶם אֲשֶׁר הֵם זֹבְחִים עַל פְּנֵי הַשָּׂדֶה וֶהֱבִיאֻם לַיהֹוָה אֶל פֶּתַח אֹהֶל מוֹעֵד אֶל הַכֹּהֵן וְזָבְחוּ זִבְחֵי שְׁלָמִים לַיהֹוָה אוֹתָם׃}
}
{
	\eP{
		So that the children of Israel will bring their
		sacrifices—which they are making at the open field—and
		will bring them to YHWH, to the entrance of the Tent of
		Meeting, to the priest, and make them sacrifices of peace
		offerings to YHWH.
	}
}
{
}

\Verse{6}
{
	\hP{וְזָרַק הַכֹּהֵן אֶת הַדָּם עַל מִזְבַּח יְהֹוָה פֶּתַח אֹהֶל מוֹעֵד וְהִקְטִיר הַחֵלֶב לְרֵיחַ נִיחֹחַ לַיהֹוָה׃}
}
{
	\eP{
		And the priest shall fling the blood on YHWH’s altar at
		the entrance of the Tent of Meeting and burn the fat to
		smoke as a pleasant smell to YHWH.
	}
}
{
}

\Verse{7}
{
	\hP{וְלֹא יִזְבְּחוּ עוֹד אֶת זִבְחֵיהֶם לַשְּׂעִירִם אֲשֶׁר הֵם זֹנִים אַחֲרֵיהֶם חֻקַּת עוֹלָם תִּהְיֶה זֹּאת לָהֶם לְדֹרֹתָם׃}
}
{
	\eP{
		And they shall not make their sacrifices anymore to the
		satyrs, after whom they are whoring. This shall be an
		eternal law to them through their generations.
	}
}
{
%	\footnote{
%		satyrs. It is not certain what these are. The term “satyrs” is used
%		because the Hebrew term has to do with goats. They are mentioned in Isa
%		13:21; 34:14 and in 2 Chr 11:15.
%	}
}

\Verse{8}
{
	\hP{וַאֲלֵהֶם תֹּאמַר אִישׁ אִישׁ מִבֵּית יִשְׂרָאֵל וּמִן הַגֵּר אֲשֶׁר יָגוּר בְּתוֹכָם אֲשֶׁר יַעֲלֶה עֹלָה אוֹ זָבַח׃}
}
{
	\eP{
		“And you shall say to them: Any man from the house of
		Israel and from the aliens who will reside among them who
		will make a burnt offering or a sacrifice
	}
}
{
}

\Verse{9}
{
	\hP{וְאֶל פֶּתַח אֹהֶל מוֹעֵד לֹא יְבִיאֶנּוּ לַעֲשׂוֹת אֹתוֹ לַיהֹוָה וְנִכְרַת הָאִישׁ הַהוּא מֵעַמָּיו׃}
}
{
	\eP{
		and will not bring it to the entrance of the Tent of
		Meeting to make it to YHWH: that man will be cut off from
		his people.
	}
}
{
}

\Verse{10}
{
	\hP{וְאִישׁ אִישׁ מִבֵּית יִשְׂרָאֵל וּמִן הַגֵּר הַגָּר בְּתוֹכָם אֲשֶׁר יֹאכַל כּל דָּם וְנָתַתִּי פָנַי בַּנֶּפֶשׁ הָאֹכֶלֶת אֶת הַדָּם וְהִכְרַתִּי אֹתָהּ מִקֶּרֶב עַמָּהּ׃}
}
{
	\eP{
		“And any man from the house of Israel and from the aliens
		who will reside among them who will eat any blood: then I
		shall set my face against the person who eats the blood,
		and I shall cut him off from among his people.
	}
}
{
%	\footnote{
%		person. This is very difficult to capture in translation. The word for
%		“person,” Hebrew nepeš, is the same word that means “life” three times
%		in v. 11. In the Hebrew, therefore, it is more obvious that the issue is
%		that a living creature is eating blood, which bears the life of another
%		living creature.
%	}
}

\Verse{11}
{
	\hP{כִּי נֶפֶשׁ הַבָּשָׂר בַּדָּם הִוא וַאֲנִי נְתַתִּיו לָכֶם עַל הַמִּזְבֵּחַ לְכַפֵּר עַל נַפְשֹׁתֵיכֶם כִּי הַדָּם הוּא בַּנֶּפֶשׁ יְכַפֵּר׃}
}
{
	\eP{
		Because the flesh’s life is in the blood, and I have given
		it to you on the altar to make atonement over your lives,
		because it is the blood that makes atonement for life.
	}
}
{
}

\Verse{12}
{
	\hP{עַל כֵּן אָמַרְתִּי לִבְנֵי יִשְׂרָאֵל כּל נֶפֶשׁ מִכֶּם לֹא תֹאכַל דָּם וְהַגֵּר הַגָּר בְּתוֹכְכֶם לֹא יֹאכַל דָּם׃}
}
{
	\eP{
		On account of this I have said to the children of Israel:
		every person among you shall not eat blood, and the alien
		who resides among you shall not eat blood.
	}
}
{
}

\Verse{13}
{
	\hP{וְאִישׁ אִישׁ מִבְּנֵי יִשְׂרָאֵל וּמִן הַגֵּר הַגָּר בְּתוֹכָם אֲשֶׁר יָצוּד צֵיד חַיָּה אוֹ עוֹף אֲשֶׁר יֵאָכֵל וְשָׁפַךְ אֶת דָּמוֹ וְכִסָּהוּ בֶּעָפָר׃}
}
{
	\eP{
		“And any man from the children of Israel and from the
		aliens who reside among them who will hunt game, animal or
		bird, that may be eaten: he shall spill out its blood and
		cover it with dust.
	}
}
{
}

\Verse{14}
{
	\hP{כִּי נֶפֶשׁ כּל בָּשָׂר דָּמוֹ בְנַפְשׁוֹ הוּא וָאֹמַר לִבְנֵי יִשְׂרָאֵל דַּם כּל בָּשָׂר לֹא תֹאכֵלוּ כִּי נֶפֶשׁ כּל בָּשָׂר דָּמוֹ הִוא כּל אֹכְלָיו יִכָּרֵת׃}
}
{
	\eP{
		Because all flesh’s life: its blood is one with its life.
		So I say to the children of Israel: you shall not eat the
		blood of all flesh—because all flesh’s life: it is its
		blood. Everyone of those who eat it will be cut off.
	}
}
{
}

\Verse{15}
{
	\hP{וְכל נֶפֶשׁ אֲשֶׁר תֹּאכַל נְבֵלָה וּטְרֵפָה בָּאֶזְרָח וּבַגֵּר וְכִבֶּס בְּגָדָיו וְרָחַץ בַּמַּיִם וְטָמֵא עַד הָעֶרֶב וְטָהֵר׃}
}
{
	\eP{
		And any living being, whether a citizen or an alien, who
		will eat carcass or a torn animal shall wash his clothes
		and shall wash with water and will be impure until evening
		and then will be pure.
	}
}
{
%	\footnote{
%		carcass or a torn animal. See the comments on Lev 7:24.
%	}
}

\Verse{16}
{
	\hP{וְאִם לֹא יְכַבֵּס וּבְשָׂרוֹ לֹא יִרְחָץ וְנָשָׂא עֲוֺנוֹ׃}
}
{
	\eP{
		And if he will not wash them and will not wash his flesh,
		then he shall bear his crime.”
	}
}
{
}
